\section{Large-Scale Events, Epidemic Contagion, and Vaccination}
\label{sec:results_large_scale_events_epidemics}

Large-scale gatherings and infectious-disease outbreaks provide two complementary tests of population simulation at the city scale. A large event produces a short, concentrated demand for movement that competes with ordinary urban routines and, when several events overlap, with other event destinations. An epidemic instead develops over many days as population movement repeatedly changes which susceptible and infected groups share the same locations. Vaccination and movement restrictions then make it possible to examine how public-health interventions modify this longer-term process.

These use cases are particularly relevant to LodusPop's environment-driven design. Event locations actively request attendees rather than requiring individual agents to maintain event preferences, while points of interest apply infection dynamics to the population currently present. The same population blobs can consequently be divided, relocated, assigned new traceable states, and later merged when their states are compatible. The experiments therefore evaluate the interaction between short-term mobility demands, long-term population-state changes, and policy actions in a common 94-region representation of Porto Alegre.

\subsection{Modeled Plugins}
\label{sec:large_scale_epidemic_plugins}

\subsubsection{Gather-Population Plugin}

The \textit{Gather-Population Plugin} introduces the \texttt{gather\_population} action used to represent sports, commercial, and cultural events. A gather action specifies a destination point of interest $P_d$, a requested population quantity, and a population template. The plugin searches points of interest containing eligible population and assigns an origin contribution according to population availability and distance to $P_d$. Under the distance-weighted configuration used in the published experiments, nearby origins receive greater weight.

After the contribution of each origin has been determined, the complex gather action is decomposed into multiple \texttt{move\_population} actions. Each sub-action transfers its assigned population from an origin point of interest to the event destination. The requested quantity is therefore a demand rather than a guarantee: the quantity actually gathered depends on eligible population available when the action is executed. This distinction becomes important when concurrent events draw from the same city population.

\subsubsection{Infection Plugin}

The \textit{Infection Plugin} implements the susceptible--infected--removed (SIR) model~\citep{kermack1927contribution}. Every blob receives a traceable \texttt{sir\_status} characteristic with one of three values: \texttt{susceptible}, \texttt{infected}, or \texttt{removed}. The removed state represents people who no longer contribute to contagion, either because they recovered and acquired immunity or because they died. Blobs with different SIR states may coexist at the same point of interest, but they remain separate population units because their traceable characteristics differ.

For a point of interest containing a total population $N$, the SIR dynamics are

\begin{equation}
\label{eq:large_scale_epidemic_delta_s}
   \frac{dS}{dt}=-\frac{\beta I S}{N},
\end{equation}

\begin{equation}
\label{eq:large_scale_epidemic_delta_i}
   \frac{dI}{dt}=\frac{\beta I S}{N}-\gamma I,
\end{equation}

\begin{equation}
\label{eq:large_scale_epidemic_delta_r}
   \frac{dR}{dt}=\gamma I,
\end{equation}

where $S$, $I$, and $R$ are the susceptible, infected, and removed populations, respectively. The parameters $\beta$ and $\gamma$ are the infection and removal rates, and the basic reproduction ratio is $R_0=\beta/\gamma$. The \texttt{infect} action evaluates these quantities for the population currently sharing each point of interest and changes the traceable SIR states while conserving the total represented population.

\subsubsection{Vaccination and Movement-Restriction Plugins}

The \textit{Vaccine Plugin} adds a traceable \texttt{vaccine\_level} characteristic and a \texttt{vaccinate} action. Vaccine level is an integer that records the number of doses received. When vaccination is requested, eligible people are gathered at the pharmacy in their current region, their vaccine level is increased, and the vaccinated blobs return to their previous locations. Later doses are restricted to people who received the preceding dose.

Each vaccine level is associated with an effectiveness value. When susceptible population would transition to the infected state through Equation~\ref{eq:large_scale_epidemic_delta_i}, the transition is multiplied by one minus the effectiveness associated with that population's vaccine level. For example, an effectiveness of 0.4 reduces the modeled susceptible-to-infected transition by 40\%. The number of vaccines administered per day follows the Porto Alegre vaccination series used in the original study and is distributed among the 94 regions in proportion to their populations.

The epidemic-policy experiments also reuse the \texttt{Restriction} parameter of the original L\'evy Walk plugin described in Section~\ref{sec:results_routine_mobility}. This factor reduces routine movement between regions. A value of 0.50 retains half of the otherwise requested movement, while a value of 0.75 represents a stronger restriction in which approximately 75\% of the affected population does not undertake the inter-region movement. The restriction changes mobility rather than directly changing the SIR equations.

\subsection{Simulation Setup}
\label{sec:large_scale_epidemic_setup}

\subsubsection{Shared 94-Region Environment}

All experiments in this section use the complete neighborhood-level representation of Porto Alegre. The environment contains 94 regions and 504 points of interest and represents a population of 1,409,351. Each region contains home, work, school, marketplace, pharmacy, and gas-station points of interest; selected regions additionally contain stadiums, shopping malls, hospitals, universities, or other specialized destinations. One population blob is initialized at the home point of interest of each region according to neighborhood population data. Its sampled characteristics represent age and occupation, while its traceable region of origin allows displaced population to return home.

A simulation cycle contains 24 hourly steps and represents one day. Both application families use the published work-and-school routine as their mobility baseline, with the original L\'evy Walk scale set to $\texttt{levy\_scale}=2$. Thus, the event experiments add occasional gathering actions to ordinary mobility, whereas the epidemic experiments repeatedly apply infection and public-health actions as workers and students move through the same full-city environment. No reduced-region scenario is included in this section.

\subsubsection{Large-Scale Event Scenarios}

Four large-scale events were configured at distinct points of interest:

\begin{itemize}
    \item Event A: a sports event at the stadium in the Farrapos region;
    \item Event B: a sports event at the stadium in the Praia de Belas region;
    \item Event C: a commercial event at the marketplace in the Centro Hist\'orico region; and
    \item Event D: a cultural event at the university in the Partenon region.
\end{itemize}

Each event requests 500,000 attendees. The experiment catalog contains a no-event baseline and every combination of one, two, three, or four simultaneous events. In the simultaneous conditions, all included events occur at cycle step 19, corresponding to 19:00. If all four events occur together, their combined request is 2,000,000 people, which exceeds the city's represented population and forces the actions to compete for population available at their execution time. A final condition retains all four events but schedules Event C at step 9 and Event D at step 14, reducing their temporal overlap with the events held at step 19. Each scenario covers one 24-step cycle.

\subsubsection{Epidemic and Vaccination Scenarios}

The epidemic experiments run for 500 cycles, or 500 simulated days. At initialization, 50 people in each of the 94 regions are assigned the infected state, producing 4,700 initially infected people. No one begins in the removed state, and the remainder of the population begins susceptible. The infection action is included in the global routine and executes at every cycle step.

Three infection-only scenarios vary $R_0$ using values reported in studies of COVID-19 transmission in Brazil~\citep{dutra2020covidR0,deSouza2020covirR0,dutra2021covidR0}. The removal rate is fixed at $\gamma=1/15$, corresponding to a 15-day transition period from infected to removed, while $\beta$ is set to 0.1733, 0.204, or 0.4. These values produce $R_0$ configurations of 2.6, 3.07, and 6.0, respectively. The $R_0=2.6$ condition is used as the reference for the vaccination and restriction experiments.

Table~\ref{tbl:vaccine_configuration_94} summarizes the four published vaccination configurations. The starting day indicates when each dose becomes available, and a later dose can be administered only to people who received the preceding dose. Configurations A--C progressively add doses and effectiveness, whereas Configuration D represents an idealized single-dose vaccine with 100\% modeled effectiveness.

\begin{table}[ht]
\centering
\footnotesize
\caption{Vaccination configurations used in the 94-region epidemic experiments. Effectiveness values and starting days are listed in dose order.}
\label{tbl:vaccine_configuration_94}
\begin{tabular}{@{}lccc@{}}
\toprule
\textbf{Configuration} & \textbf{Dose count} & \textbf{Effectiveness} & \textbf{Starting day} \\
\midrule
Vaccination A & 1 & [0.4] & [50] \\
Vaccination B & 2 & [0.4, 0.7] & [50, 150] \\
Vaccination C & 3 & [0.4, 0.7, 1.0] & [50, 150, 250] \\
Vaccination D & 1 & [1.0] & [50] \\
\bottomrule
\end{tabular}
\end{table}

The final policy matrix combines no restriction, a 0.50 restriction, or a 0.75 restriction with either no vaccination or Vaccination C. Vaccination C is used instead of idealized Vaccination D because it represents increasing protection over a multi-dose schedule. This design separates the effect of mobility reduction, vaccination, and their combination under the same $R_0=2.6$ infection setting.

\subsection{Results}
\label{sec:large_scale_epidemic_results}

\subsubsection{Large-Scale Event Results}

Table~\ref{tbl:large_scale_event_displacements_94} reports total displacements, average displacements per inhabitant, and the increase relative to ordinary mobility. A single event increases the total displacement count by between 27.64\% and 28.40\%, depending on its location and the population available nearby. The differences among single-event conditions are small in aggregate, although their distance distributions differ because each destination draws from a different part of the city.

\begin{table}[ht]
\footnotesize
\centering
\caption{Results of the 94-region large-scale event simulations. A scenario name lists the events added to the ordinary-mobility baseline.}
\label{tbl:large_scale_event_displacements_94}
\begin{tabular}{@{}lrrr@{}}
\toprule
\textbf{Scenario} & \shortstack{\textbf{Total}\\\textbf{displacements}} & \shortstack{\textbf{Displacements}\\\textbf{per inhabitant}} & \shortstack{\textbf{Increase over}\\\textbf{baseline}} \\
\midrule
Baseline               & 2,570,658 & 1.8240 & --- \\
\midrule
Baseline+A             & 3,290,469 & 2.3347 & 28.00\% \\
Baseline+B             & 3,300,782 & 2.3421 & 28.40\% \\
Baseline+C             & 3,289,002 & 2.3337 & 27.94\% \\
Baseline+D             & 3,281,102 & 2.3281 & 27.64\% \\
\midrule
Baseline+AB            & 3,786,856 & 2.6870 & 47.31\% \\
Baseline+AC            & 3,748,417 & 2.6597 & 45.82\% \\
Baseline+AD            & 3,804,017 & 2.6991 & 47.98\% \\
Baseline+BC            & 3,738,497 & 2.6526 & 45.43\% \\
Baseline+BD            & 3,839,126 & 2.7240 & 49.34\% \\
Baseline+CD            & 3,826,094 & 2.7148 & 48.84\% \\
\midrule
Baseline+ABC           & 4,042,413 & 2.8683 & 57.25\% \\
Baseline+ABD           & 4,109,929 & 2.9162 & 59.88\% \\
Baseline+ACD           & 4,069,223 & 2.8873 & 58.29\% \\
Baseline+BCD           & 4,081,563 & 2.8961 & 58.78\% \\
\midrule
Baseline+ABCD          & 4,235,926 & 3.0056 & 64.78\% \\
Baseline+ABCD+DiffStep & 4,638,329 & 3.2911 & 80.43\% \\
\bottomrule
\end{tabular}
\end{table}

Figure~\ref{fig:large_scale_event_displacements_94} presents the corresponding displacement-distance distributions. The first panel compares the four single-event conditions. The second compares scenarios containing three simultaneous events. The final panel compares four simultaneous events with the staggered four-event condition.

\begin{figure}[!htbp]
  \centering
  \subfloat[Baseline and single-event scenarios.]
  {\includegraphics[width=0.95\linewidth]{Figures/Displacement_SingleEvent.png}
  \label{fig:large_scale_event_single_94}}\\
  \subfloat[Baseline and three-event scenarios.]
  {\includegraphics[width=0.95\linewidth]{Figures/Displacement_ThreeEvents.png}
  \label{fig:large_scale_event_three_94}}\\
  \subfloat[Baseline, four simultaneous events, and four events with different starting steps.]
  {\includegraphics[width=0.95\linewidth]{Figures/Displacement_AllEvents.png}
  \label{fig:large_scale_event_all_94}}
  \caption{Population-displacement distributions for the 94-region large-scale event scenarios.}
  \label{fig:large_scale_event_displacements_94}
\end{figure}

The growth in displacement is sub-additive when events occur together. For example, Events A, B, and C individually increase displacement by 28.00\%, 28.40\%, and 27.94\%, respectively, but the simultaneous ABC scenario increases displacement by 57.25\%, rather than by the sum of the three independent increases. Part of this difference occurs because people can travel directly from work or school to an event without first returning home. More importantly, simultaneous gather actions compete for the population available in nearby regions, reducing attendance that later event requests can obtain.

Temporal separation changes that result. Four simultaneous events produce 4,235,926 displacements, 64.78\% above the baseline. When Events C and D occur earlier in the day, the total rises to 4,638,329, or 80.43\% above baseline. Staggering reduces direct competition at a single step and allows some people to participate in more than one event during the day, increasing their individual displacement counts.

\subsubsection{Epidemic-Contagion Results}

Figure~\ref{fig:epidemic_r0_94} shows cumulative infections under the three reproduction-ratio settings. The curves include the 4,700 people infected at initialization and every subsequent transition from susceptible to infected. Increasing $R_0$ accelerates contagion and increases the cumulative infected population over the simulated period.

\begin{figure}[ht]
 \centering
 \includegraphics[width=0.95\linewidth]{Figures/InfectedR0.png}
 \caption{Cumulative infected population over 500 days for the three $R_0$ configurations in the 94-region environment.}
 \label{fig:epidemic_r0_94}
\end{figure}

The $R_0=2.6$ reference reaches a cumulative total of 875,962 infected people. Cumulative incidence should not be confused with the number simultaneously infected: as Figure~\ref{fig:epidemic_sir_status_94} shows, the active infected population peaks at 81,528 people on day 247, while earlier infections progressively enter the removed state. The result demonstrates that traceable state changes remain attached to mobile population blobs across hundreds of daily cycles and can be aggregated into both cumulative and point-in-time measures.

\begin{figure}[ht]
 \centering
 \includegraphics[width=0.95\linewidth]{Figures/SIRPopulationStatus.png}
 \caption{Susceptible, infected, and removed populations over time for the $R_0=2.6$ scenario.}
 \label{fig:epidemic_sir_status_94}
\end{figure}

\subsubsection{Vaccination and Movement-Restriction Results}

Vaccination substantially reduces cumulative infections relative to the $R_0=2.6$ reference. Table~\ref{tbl:vaccination_results_94} collects the published final totals and their reduction relative to the unvaccinated reference. Even Vaccination A, whose single dose has an effectiveness of 0.4, reduces the cumulative total from 875,962 to 440,569. Adding a second dose in Vaccination B lowers the result to 360,607. Vaccination C reaches 355,124 infections, only modestly below Vaccination B because its third, fully effective dose does not begin until day 250. The idealized Vaccination D produces the lowest vaccination-only total, 170,978.

\begin{table}[ht]
\centering
\footnotesize
\caption{Final cumulative infected population in the published vaccination scenarios. Reductions are calculated relative to the unvaccinated $R_0=2.6$ reference.}
\label{tbl:vaccination_results_94}
\begin{tabular}{@{}lrr@{}}
\toprule
\textbf{Scenario} & \textbf{Cumulative infected} & \textbf{Reduction} \\
\midrule
Unvaccinated reference & 875,962 & --- \\
Vaccination A & 440,569 & 49.70\% \\
Vaccination B & 360,607 & 58.83\% \\
Vaccination C & 355,124 & 59.46\% \\
Vaccination D & 170,978 & 80.48\% \\
\bottomrule
\end{tabular}
\end{table}

\begin{figure}[ht]
 \centering
 \includegraphics[width=0.95\linewidth]{Figures/InfectedVaccines.png}
 \caption{Cumulative infected population for the vaccination configurations in Table~\ref{tbl:vaccine_configuration_94}. All scenarios use $R_0=2.6$.}
 \label{fig:vaccination_results_94}
\end{figure}

Movement restriction alone also reduces cumulative infections compared with unrestricted mobility. Across the combined policy matrix, however, Vaccination C produces a larger reduction than movement restriction alone, and the most favorable outcomes occur when both policies are active. The condition combining Vaccination C with a restriction of 0.75 reaches 140,941 cumulative infections, an 83.91\% reduction from the unvaccinated, unrestricted reference. This total is also lower than the 170,978 infections obtained with the idealized Vaccination D alone, illustrating the additional effect produced by combining pharmaceutical and mobility interventions under the modeled assumptions.

\begin{figure}[ht]
 \centering
 \includegraphics[width=0.95\linewidth]{Figures/InfectedVaccineRestriction.png}
 \caption{Epidemic scenarios combining no restriction, a 0.50 restriction, or a 0.75 restriction with either no vaccination or Vaccination C. All scenarios use $R_0=2.6$.}
 \label{fig:vaccination_restriction_results_94}
\end{figure}

\subsection{Remarks}

The large-event experiments show how an environment-driven request can produce citywide movement without assigning a detailed event-choice model to every individual. Their central result is not simply that events add trips, but that the result depends on competition and timing. Concurrent events cannot be analyzed as independent additions to baseline mobility because they draw from the same finite population. Staggered events reduce that competition and allow repeated participation, producing more total displacement than the simultaneous four-event condition.

The epidemic experiments extend the same population representation from hourly movement to state evolution over 500 days. SIR status and vaccine level demonstrate why traceable characteristics are useful: population groups can split when their health state changes, preserve that state as they move, and interact with other groups at subsequent destinations. The vaccination action further demonstrates composition between environmental requests, mobility, and population-state transitions, while the restriction factor modifies contact opportunities indirectly by reducing routine movement.

The numerical outcomes should be interpreted as comparisons among the implemented scenarios rather than forecasts for Porto Alegre. The SIR process uses homogeneous $\beta$ and $\gamma$ parameters within each configuration, vaccine effectiveness acts directly on modeled infection transitions, daily vaccination demand follows the input series, and event attendance is driven by availability and distance rather than individual preferences. Within those assumptions, the experiments reproduce the published findings: event timing changes competition for population, higher $R_0$ accelerates epidemic growth, vaccination reduces cumulative infection, and vaccination combined with movement restriction produces the lowest infection total among the reported policy scenarios.
